\pdfoutput=1
\documentclass[11pt,a4paper]{article}
\usepackage{jheppub}
\usepackage{amsthm}
\usepackage{wrapfig}
\usepackage{pgfplots}
\usepackage[makeroom]{cancel}
\pgfplotsset{compat=1.18}
\notoc

\title{Introduction to Supersymmetry}
\author{Edward Witten\note{Supported in part by NSF Grant PHY80-19754.}}
\affiliation{Joseph Henry Laboratories,\\
Princeton University,\\
Princeton, New Jersey 08544}

\begin{document}
\maketitle
\flushbottom

\section{Introduction}

Supersymmetry is a remarkable subject that has fascinated
particle physicists since it was originally introduced.\textsuperscript{1}
Although supersymmetry is no longer a new idea, we still do not know in what
form, if any, it plays a role in the proper description of nature.

The present status of supersymmetry might be compared very
roughly to the status of non-Abelian gauge theories twenty years ago.
It is a fascinating mathematical structure, and a reasonable extension
of current ideas, but plagued with phenomenological difficulties.

In these lectures, I will present an introduction to supersymmetry,
or at least to some aspects of this extensive subject. I
also will describe some recent results. Supersymmetry has the reputation
of a subject that is difficult to learn. I will try to at
least partially dispel this unjustified impression.

In these lectures we will discuss only global supersymmetry,
not supergravity.

A number of useful reviews of supersymmetry are available,
which approach the subject from several points of view.\textsuperscript{2}
Many aspects of supersymmetry that I will not be able to discuss are
treated in these review articles.

\section{Bose-Fermi Symmetry}

If one considers a theory of two decoupled Bose fields, $\phi_1$, and
$\phi_2$, so
\begin{equation}
\mathcal{L} = \int d^4x \, \frac{1}{2}\left((\partial_\mu \phi_1)^2 + (\partial_\mu \phi_2)^2\right)
\tag{1}
\end{equation}
it is possible to combine $\phi_1$ and $\phi_2$ into a conserved current
\[
J_\mu = \phi_1\, \partial_\mu \phi_2 - \phi_2\, \partial_\mu \phi_1
\]
even though they are not interacting.
Actually, since, for example, $\partial_\alpha \phi_1$ or
$\partial_\alpha \partial_\beta \phi_1$ satisfies the Klein-Gordon
equation just as $\phi_1$ does, one can form in the free field
theory additional conserved currents such as
\begin{align}
J_{\mu\alpha} &= \partial_\alpha \phi_1\, \partial_\mu \phi_2 - \phi_2\, \partial_\mu \partial_\alpha \phi_1 \notag\\
J_{\mu\alpha\beta} &= \partial_\alpha \partial_\beta \phi_1\, \partial_\mu \phi_2 - \phi_2\, \partial_\mu \partial_\alpha \partial_\beta \phi_1
\tag{2}
\end{align}
It is easy to check that these are conserved because
$\nabla^2 \phi_1 = \nabla^2 \phi_2 = 0$.

The reason that $J_\mu$ is far more important than $J_{\mu\alpha}$ or
$J_{\mu\alpha\beta}$ is that $J_\mu$ can still be conserved in interacting
theories. One can add to $\mathcal{L}$ a term like $V(\phi_1^2 + \phi_2^2)$
that is invariant under $\delta \phi_1 = \phi_2$, $\delta \phi_2 = -\phi_1$.
$J_\mu$ is still conserved.

However, when interactions are included, $J_{\mu\alpha}$ and $J_{\mu\alpha\beta}$
are no longer conserved. Moreover, it is not possible to redefine them (by
adding extra terms to allow for the interactions) so that they will still be
conserved. One may readily verify this in special cases. In general, it is a
consequence of the Coleman-Mandula theorem.\textsuperscript{3} Coleman and
Mandula showed (basically by S matrix theory alone) that in a theory with
non-zero scattering amplitudes in more than 1+1

% Original pages: 307--311
% Figures: Figure 1 placeholder on original p. 308.
% Uncertainty log: none after manual scan/OCR cross-check.

dimensions the only possible conserved quantities that transform as
tensors under the Lorentz group are the following. The usual space-
time symmetries are certainly allowed: the energy-momentum operator
$P_\mu$ and the Lorentz transformations $M_{\alpha\beta}$ commute with the
$S$ matrix in all of our usual theories. We may also have arbitrary Lorentz
invariant conserved quantum numbers $Q_i$ (electric charge, baryon
number, etcs.). Finally, if all particles are massless, the Coleman-
Mandula theorem allows conformal invariance, which is not usually
realized, however, in field theories with interactions. The Coleman-
Mandula theorem forbids ``exotic'' conservation laws --- conservation
laws other than the usual space-time symmetries which do not commute
with Lorentz transformations.

For a proof, I refer you to reference (3); I will just give a
rough idea. The basic idea is that conservation of $P_\mu$ and $M_{\alpha\beta}$
leaves only the scattering angle unknown in (say) a two body
collision. Additional, exotic conservation laws would determine the
scattering angle, leaving only a discrete set of possible angles.
Since the scattering amplitude is always an analytic function of
angle, it would actually then have to vanish at all angles.

Note that the argument obviously does not apply in $1+1$
dimensions. In $1+1$ dimensions, the only possible angles are $0$ and $\pi$;
there is no such thing as analyticity as a function of scattering
angle. This is why, in $1+1$ dimensions, it is possible to have
interacting systems, such as the sine-gordon equation, with exotic
conservation laws.

To illustrate the argument by a concrete example, suppose we
have a conserved traceless symmetric tensor $Q_{\beta\gamma}$. This is an exotic
conservation law because, transforming as a tensor, it does not
commute with Lorentz transformations. Its matrix element in a one
particle state of momentum $p$ and (for simplicity) spin zero would
have to be $\langle p| Q_{\beta\gamma} |p\rangle = p_\beta p_\gamma - \frac{1}{4}\delta_{\beta\gamma} p^2$, by Lorentz invariance.

Applied to a two-body collision, with incident particles of momentum
$p_1$, $p_2$, and outgoing particles of momentum $q_1$, $q_2$ (figure (1)), the
conservation of $Q_{\beta\gamma}$ would tell us $p_{1\beta}\, p_{1\gamma} + p_{2\beta}\, p_{2\gamma} = q_{1\beta}\, q_{1\gamma} +
q_{2\beta}\, q_{2\gamma}$. This is possible only if the scattering angle is zero.\begingroup\renewcommand{\thefootnote}{\fnsymbol{footnote}}\footnotemark[1]\endgroup
With more effort, the same type of argument works even if the
particles have non-zero spin.

Now going back to the exotic conserved currents that we defined
in free field theory (equation (2)), the corresponding conserved
charges
\begin{equation}
\begin{aligned}
Q_\alpha &= \int d^3x\, J_{0\alpha} \\
Q_{\alpha\beta} &= \int d^3x\, J_{0\alpha\beta}
\end{aligned}
\tag{3}
\end{equation}
are forbidden by the Coleman-Mandula theorem for theories with a
non-trivial $S$ matrix in more than $1+1$ dimension. This is why it is
impossible to add interactions to
the free field theory in a way that
preserves the conservation of $J_{\mu\alpha}$
and $J_{\mu\alpha\beta}$.

% TODO FIGURE: Original Fig. 1 appears here on p. 308.
\begin{wrapfigure}{l}{0.28\textwidth}
\vspace{-1.3\baselineskip}
\refstepcounter{figure}
\centering
\fbox{\parbox[c][1.25in][c]{0.22\textwidth}{\centering TODO\\Figure placeholder\\Original p.~308}}
\par\smallskip
Figure~\thefigure
\vspace{-0.9\baselineskip}
\end{wrapfigure}

\begingroup
\renewcommand{\thefootnote}{\fnsymbol{footnote}}
\footnotetext[1]{It is here assumed that the matrix element of $Q_{\beta\gamma}$ in the two particle state $|p_1 p_2\rangle$ is the sum of the matrix elements in the states $|p_1\rangle$ and $|p_2\rangle$. This is true if $Q_{\beta\gamma}$ is the integral of a local current density, or, more generally, if $Q_{\beta\gamma}$ is defined in a way that is ``not too non-local''.}
\endgroup

Instead of two non-interacting scalars, consider next a free
massless charged scalar $\phi$ and a free two component (left handed)
fermion $\psi_\alpha$. The Lagrangian is
\begin{equation}
\mathcal{L} = \int d^4x\, \left(\partial_\mu \phi^* \partial^\mu \phi + \bar\psi\, i \not\!\partial\, \psi \right)
\tag{4}
\end{equation}
Again --- although the fields are non-interacting --- one can define
conserved currents connecting them. One of the simplest is\begingroup\renewcommand{\thefootnote}{\fnsymbol{footnote}}\footnotemark[1]\endgroup
\begin{equation}
S_{\mu\alpha} = \left(\partial_\sigma \phi^* \gamma^\sigma \gamma_\mu \psi\right)_\alpha
\tag{5}
\end{equation}
It is easy to show $\partial_\mu S^\mu_{\ \alpha} = 0$:
\begin{equation}
\partial_\mu S^\mu_{\ \alpha} = \left(\partial_\mu \partial_\sigma \phi^* \gamma^\sigma \gamma^\mu \psi\right)_\alpha + \left(\partial_\sigma \phi^* \gamma^\sigma \gamma^\mu \partial_\mu \psi\right)_\alpha
\tag{6}
\end{equation}
The second term vanishes since $\gamma^\mu \partial_\mu \psi = 0$. For the first term, note
$\partial_\mu \partial_\sigma \phi$ is symmetric in $\mu$ and $\sigma$, so we may replace $\gamma^\sigma \gamma^\mu$ by
$1/2(\gamma^\mu \gamma^\sigma + \gamma^\sigma \gamma^\mu) = g^{\mu\sigma}$. We then have $g^{\mu\sigma}\partial_\mu \partial_\sigma \phi^* = 0$.

Again, additional conserved currents can be defined. The only
property of $\psi$ that we used was the Dirac equation, which is also
satisfied by $\partial_\gamma \psi$, so
\begin{equation}
S_{\mu\gamma\alpha} = \left(\partial_\sigma \phi^* \gamma^\sigma \gamma_\mu \partial_\gamma \psi\right)_\alpha
\tag{7}
\end{equation}
is also conserved, that is, $\partial_\mu S^\mu_{\ \gamma\alpha} = 0$.

The exciting fact that makes supersymmetry interesting is that,
although conservation of $S_{\mu\gamma\alpha}$ is always ruined when interactions are
included, it is possible to add interactions in such a way that $S_{\mu\alpha}$
is conserved. A simple example is
\begin{equation}
\mathcal{L} = \int d^4x\, \left(\partial_\mu \phi^* \partial^\mu \phi + \bar\psi\, i \not\!\partial\, \psi - g^2 |\phi|^4 - g(\phi^* \psi_\alpha \psi^\alpha + h.c.)\right)
\tag{8}
\end{equation}
which is known as the ``massless Wess-Zumino model''. Although $S_{\mu\alpha}$
as previously defined is not conserved in this model, by adding an
extra term
\begin{equation}
S_{\mu\alpha} \rightarrow S_{\mu\alpha} + i g\, \gamma_\mu \phi^{*2}\, \psi^*
\tag{9}
\end{equation}
One preserves conservation of the current in the presence of
interaction. (Here $\psi^*$ is a right-handed two component spinor, the
hermitian conjugate of $\psi$, which is left-handed. See, for example, the
notes of Wess and Bagger\textsuperscript{2} for more details.)

\begingroup
\renewcommand{\thefootnote}{\fnsymbol{footnote}}
\footnotetext[1]{$S_{\mu\alpha}$ is here a left-handed Weyl spinor, like $\psi$ itself. We will soon go over to a Majorana basis.}
\endgroup

Why can $S_{\mu\alpha}$, but not $S_{\mu\gamma\alpha}$, be conserved in the presence of
interactions? We can find out by studying the conserved charges
\begin{equation}
\begin{aligned}
\hat Q_\alpha &= \int d^3x\, S_{0\alpha} \\
\hat Q_{\gamma\alpha} &= \int d^3x\, S_{0\gamma\alpha}
\end{aligned}
\tag{10}
\end{equation}
We cannot apply the Coleman-Mandula theorem directly to the $\hat Q_\alpha$ and
$\hat Q_{\gamma\alpha}$ because the Coleman-Mandula theorem deals with conserved charges
that transform as Lorentz tensors, while the $\hat Q$ transform as spinors
(or vector-spinors). However, the Coleman-Mandula theorem can be
applied to the bosonic conserved charges that can be formed from the
anticommutators of the $\hat Q$. (It is the anti-commutators of the $\hat Q$ that
we should consider, because $S_{\mu\alpha}$ and $S_{\mu\gamma\alpha}$ being linear in fermi
fields, anti-commute at spacelike separation.)

Now $\hat Q_\alpha$, being a left-handed spinor, transforms as $(1/2,0)$
under Lorentz transformations. Its hermitian adjoint, $\hat Q^*_{\dot\alpha}$, transforms
as $(0,1/2)$. The anticommutator of $\hat Q$ with its adjoint $\hat Q^*$, which cannot
vanish (since the anticommutator of any operator with its hermitian adjoint is non-zero), transforms as $(1/2,1/2)$ under Lorentz
transformations. The Coleman-Mandula theorem permits the conservation
of precisely one operator that transforms as $(1/2,1/2)$, namely
the energy-momentum operator $P_\mu$. But $\hat Q_{\gamma\alpha}$, a vector-spinor, has
components of spin up to $3/2$. The anti-commutator $\left\{\hat Q_{\gamma\alpha}, \hat Q^*_{\sigma\tau}\right\}$ which
cannot vanish, for the reason noted above, has components of spin up
to three. Since $\left\{\hat Q_{\gamma\alpha}, \hat Q^*_{\sigma\tau}\right\}$ is conserved if $\hat Q_{\gamma\alpha}$ is, and since the
Coleman-Mandula theorem does not permit conservation of an operator
of spin three in an interacting theory, $\hat Q_{\gamma\alpha}$ cannot be conserved in an
interacting theory.

The Coleman-Mandula theorem does permit conservation of the $\hat Q_\alpha$.
Changing basis from the $\hat Q_\alpha$ (not hermitian as I have defined them) to
a basis of four hermitian operators $Q_\alpha$ that transform as a (real,
Majorana) Lorentz four-spinor, the algebra of the $Q_\alpha$ turns out to be
\begin{equation}
\{Q_\alpha, \overline{Q}_\beta\} = \gamma^\mu_{\alpha\beta} P_\mu
\tag{11}
\end{equation}
One may readily check that this is so in, for example, the free
field theory (4).

The right-hand side of (11) certainly contains only operators
that are permitted by the Coleman-Mandula theorem, but one might ask
whether there are more general possibilities. This question was
answered by Haag, Sohnius, and Lopuszanski.\textsuperscript{4} Since $Q_\alpha$ transforms as
$(1/2,0) + (0,1/2)$, its anti-commutator with itself might contain
pieces transforming as $(1,0)$, $(0,1)$, or $(0,0)$, apart from the
$(1/2,1/2)$ term we have already encountered. By the Coleman-Mandula
theorem, the only conserved operators transforming as $(1,0)$ or $(0,1)$
are the Lorentz generators $M_{\mu\nu}$; however, Haag, Sohnius, and Lopuszanski
showed that it is impossible to introduce $M_{\mu\nu}$ in (11) without
violating the Jacobi identity. As for the possibility of adding to
(11) operators that transform as $(0,0)$ (in other words, operators
that commute with Lorentz transformations), this is possible, but
only in the ``extended supersymmetry'' theories in which there are
several conserved spinors $Q_{\alpha i}$. Beautiful though those theories are,
I believe that they are too restrictive to describe the physics we
know at energies much less than the Planck mass.\begingroup\renewcommand{\thefootnote}{\fnsymbol{footnote}}\footnotemark[1]\endgroup\ At such high

\begingroup
\renewcommand{\thefootnote}{\fnsymbol{footnote}}
\footnotetext[1]{The basic problem is that they require the charged fermions to form a ``real representation'' of the gauge group.}
\endgroup


\end{document}
